%!TEX root = ED_Vetor_Ponteiro.tex

\section{Estruturas de Dados Lineares}
\begin{frame}[label=notas]
	\begin{exampleblock}{Estruturas de Dados e Algoritmos}
		Temas fundamentais da Ciência da Computação, utilizados em diversas áreas do conhecimento e com diferentes propósitos de aplicação
		\begin{itemize}
			\item Estrutura de Dados: modo particular de armazenar/organizar dados 
			para facilitar acesso e modificações
			\item Algoritmo: sequência finita de instruções (para resolver uma tarefa)
		\end{itemize}
	\end{exampleblock}
	\pause\vfill
	Problema: descobrir o telefone de um conhecido.
	\begin{itemize}
		\item Algoritmo: pesquisa sequencial por nome (\alert<2>{chave de ordenação})	
		\item Estrutura de dados: lista telefônica - {[nome] : [número de telefone]}
		\pause
		\item E se quiséssemos pesquisar pelo peso?
	\end{itemize}
	\pause\vfill
	\begin{beamerboxesrounded}{}
		Uma estrutura de dados adequada pode aumentar a eficiência do algoritmo
	\end{beamerboxesrounded}
\end{frame}

\subsection{Introdução}
\begin{frame}
	Exemplos de problemas:
	\begin{enumerate}
		\item Manipular arquivos com informações de alunos. \pause
			\begin{itemize}
				\item solução: ordenar por número de matrícula
				\item operações possíveis: inserir/remover, pesquisar, etc.
				\item qual ED seria adequada? \pause \alert<3>{Lista, árvore}
			\end{itemize}
		\pause
		\item Corrigir provas. \pause
			\begin{itemize}
				\item solução: analisar cada uma individualmente
				\item operações possíveis: colocar/retirar prova
				\item qual ED seria adequada? \pause \alert<6>{Pilha}
			\end{itemize}
		\pause
		\item Lidar com alunos que querem revisar a prova. \pause
			\begin{itemize}
				\item solução: ordenar por ordem de chegada
				\item operações possíveis: entrar (ao fim)/sair (do começo)
				\item qual ED seria adequada? \pause \alert<9>{Fila}
			\end{itemize}
	\end{enumerate}
\end{frame}

\subsection{Tipos de Dados}
\begin{frame}[label=notas]
	Termos:
	\begin{itemize}
		\item Tipo de Dado: associado a um método de interpretar um padrão de bits (linguagem de programação)
			\subitem{boolean (bit - 0/1)}
		\item Tipo Abstrato de Dado (TAD): modelo de dado e operações definidas sobre ele %(tipo + operações validas) é uma metodologia de programação que tem como proposta reduzir a informação necessária para a criação/programação de um algoritmo através da abstração das variáveis envolvidas em uma única entidade fechada. A abstração de informações através do TAD permitiu a melhor compreensão dos algoritmos, maior facilidade de programação, e por consequência, aumentou a complexidade dos programas.
			\subitem{$\mathbb{N}$ e \{+, -, /, *\}}
		\item Estrutura de Dados (ED): forma concreta de se implementar um TAD \pause $\rightarrow$ \alert<2>{representação computacional do modelo matemático}
	\end{itemize}
	\pause\vfill
	Uma ED pode ser:
	\begin{itemize}
		\item homogênea: conjuntos de dados do mesmo tipo
		\item heterogênea: conjuntos de dados de diferentes tipos
	\end{itemize}
\end{frame}

\begin{frame}[label=notas]
	\framesubtitle{Primitivos}

	Dependem da linguagem de programação usada
	\uncover<2>{
		\begin{columns}
			\column{.5\textwidth}
			\column{.4\textwidth}
				\begin{tabular}{|l|l|}
				             	\multicolumn{2}{c}{\emph{Integer}} \\ \hline
				             	\textbf{Tamanho} & \textbf{Intervalo} \\ \hline
				             	\ 4 bits & [-8, 7] \\ \hline
				             	\ 8 bits & [-128, 127] \\ \hline
				             	16 bits & [-32768, 32767] \\ \hline
				             	$\cdots$ & $\cdots$ \\ \hline
				\end{tabular}	
		\end{columns}
	}
	\vspace{-6\baselineskip}
	\begin{itemize}
		\item Boolean: valor lógico
			\subitem{$b\leftarrow falso$}
		\pause
		\item Integer: valor numérico ($\in \mathbb{Z}$)
			\subitem{$i\leftarrow 2010$}
		\pause
		\item Float: valor numérico ($\in \mathbb{R}$)
			\subitem{$f\leftarrow3.14$}
		\pause
		\item Character: letras, dígitos, símbolos
			\subitem{$c\leftarrow$'@'}
		\pause
		\item String: sequência de caracteres
			\subitem{$s\leftarrow$ ``$\#2$: Gênio é $1\%$ inspiração e $99\%$ de transpiração (T. Edison)''}
		\pause
		\item \emph{\alert{Ponteiro}}
	\end{itemize}
\end{frame}

\begin{frame}[fragile,label=notas]
	\framesubtitle{Compostos}

	Criados a partir dos tipos primitivos
	\begin{itemize}
		\item Registro: 2 ou mais primitivos\\
			\begin{center}
				\begin{minipage}{.4\textwidth}
					\begin{lstlisting}
						struct numero_complexo_t{
						  float re;
						  float im;
						};
					\end{lstlisting}
				\end{minipage}
			\end{center}
    		\pause\vfill
		\item União: diferentes tipos primitivos em um único espaço de memória
			\begin{center}
				\begin{minipage}{.25\textwidth}
					\begin{lstlisting}
						union valor_u {
						  int iVal
						  float fVal
						  char cVal
						};
					\end{lstlisting}
				\end{minipage}
			\end{center}
			\begin{itemize}
				\item Simples: o programador precisa saber qual o tipo está usando
				\item Variante: o computador sabe o tipo utilizado
			\end{itemize}
	\end{itemize}
\end{frame}
\lstset{morekeywords={numero_complexo_t,valor_u}}

\begin{frame}[label=notas]
	\framesubtitle{Abstratos}

	Modelo matemático para EDs similares
	\subitem{Definido (indiretamente) pelas operações que pode realizar}
	\begin{beamerboxesrounded}{}
		Forma de organizar e relacionar dados na memória do sistema, de modo a permitir certas operações sobre esses dados, independentemente do modo como essas estruturas são implementadas.
	\end{beamerboxesrounded}
	\pause\vfill
	\begin{columns}
		\column{.3\textwidth}
			Exemplo: Pilha
			\begin{itemize}
				\item Push (inserir)
				\item Pop (remover)
			\end{itemize}
		\column{.65\textwidth}
			\includegraphics[width=.6\textwidth]{stack}
			%source: http://commons.wikimedia.org/wiki/File:Stack-svg.png	
	\end{columns}
\end{frame}